\section[A reciprocal contact completion for a singular angle]{A reciprocal contact completion for a singular angle}
\label{sec:ui-reciprocal-contact}

The following exact algebraic lemma isolates a useful mechanism. It
does not assert a complete C2 all-tests theorem.

\begin{lemma}[Reciprocal contact completion]
Let $A_0,A_2$ be symmetric matrices over a real closed field, let
$q_0$ be a column vector, and consider
\[
 f(q,V,W)=q^T(A_0+V^2A_2)q+2VW b^Tq+(c_0+c_2V^2)W^2.
\]
Suppose $(q,V,W)=(q_0,1,1)$ is a zero with all first derivatives
zero. Then
\[
 b=-(A_0+A_2)q_0,\qquad
 c_0=q_0^TA_2q_0,\qquad c_2=q_0^TA_0q_0,
\]
and the entire polynomial has the exact decomposition
\[
 \boxed{f=(q-Vq_0W)^TA_0(q-Vq_0W)
             +(Vq-q_0W)^TA_2(Vq-q_0W).}             \tag{RC1}
\]
Consequently $A_0,A_2\succeq0$ imply a sum of squares. If
$A_0\succ0$ and $q_0^TA_2q_0>0$, its nonzero zero states have
$V=\pm1$ and $q=Vq_0W$.
\end{lemma}
\begin{proof}
The $q$ equation gives the displayed value of $b$. The $V$
equation then gives $c_2=q_0^TA_0q_0$, and the $W$ equation gives
$c_0=q_0^TA_2q_0$. Expanding the right side of (RC1) proves the
identity. The zero statement follows first from the strictly positive
$A_0$ term and then from $(V^2-1)^2W^2q_0^TA_2q_0$.
\end{proof}

\paragraph{Homogeneous form without a mixed inverse.}
For a positive pure-coefficient $k$, substitute $q=Q/k$ and
$V=v/k$, and multiply (RC1) by $k^2$. It becomes
\[
 (Q-vq_0W)^TA_0(Q-vq_0W)
 +\frac1{k^2}(vQ-k^2q_0W)^TA_2(vQ-k^2q_0W).       \tag{RC2}
\]
Thus its roots use only the pure-coefficient factor $1/k$.
There is no inverse or radical of a mixed tensor element and no
comparison between a testing rate and $k$.

\paragraph{The bounded-source C2 UI leading model.}
Here $u_c=a$ is a positive bounded unit and $v_c=k$ is
infinitesimal, with $d=1$. To state a residue model, first restrict
to the subdomain where the original $n,m,\alpha$ are bounded and
the source parameters $r,g,1-a^2b$ have nonzero residues. In the
following formulas $a,b,r$ denote those residues, $z=a^2$, and
\[
 A_c=1+zr,\quad g=(1+z)b-1+r,\quad
 \Delta=1+(1+z)b+z(1+z)b^2+2z(1+z)br+z^2r^2+2zr.
\]
The source formulas give
\[
\begin{gathered}
 n=\frac{(b+r)A_c^2-r\Delta}{2r\Delta},\qquad
 m=\frac{(1-r)A_c^2}{2g\Delta},\qquad
 \alpha=\frac{z^2(1+z)b(b+r)^2}{2(1-zb)\Delta},\\
 H_0=\overline{k^2CD/(C+D)}
       =\frac{z(1+z)(b+r)b}{2\Delta},\qquad
 A_0^{\rm weak}=\overline{A/k^2}
       =\frac{A_c^2(b+r)}{2zb\Delta}.
\end{gathered}                                                   \tag{RC3}
\]
Normalize the actual contact by $W=1$. With $v=kV$, $X=k\xi$,
$Z=k\zeta$ and the leading constraint $Y=aW$, the leading
polynomial of $F/k^2$ is
\begin{align*}
 f={}&4n(1+a^2)\xi^2
 +4m\{V(1+a^2)W-\xi-a\zeta\}^2
 +4H_0V^2(a\xi-\zeta)^2\\
 &+4A_0^{\rm weak}(1+a^2)a^2W^2
 +(4\alpha+2)(1+a^2)(VaW-\zeta)^2\\
 &+2(\xi+a\zeta)^2-2(1+a^2)\zeta^2.               \tag{RC4}
\end{align*}
It has the form of the lemma, with $q=(\xi,\zeta)^T$,
\[
 q_0=(-r/b,1/(ab))^T,
\]
\[
 q^TA_0q=4n(1+a^2)\xi^2+(4m+2)(\xi+a\zeta)^2
                      +4\alpha(1+a^2)\zeta^2,
 \qquad q^TA_2q=4H_0(a\xi-\zeta)^2.
\]
Here $A_0\succ0$ and $H_0>0$. Put
$s_0=aq_{0,1}-q_{0,2}=-(1+a^2r)/(ab)\ne0$.
The exact source identities (RC3) verify all three coefficient
relations in the lemma. Hence the apparently problematic relative
normal and scale difference combine into the single root
\[
             \frac{v(aX-Z)-s_0k^2W}{k}.             \tag{RC5}
\]
The ordinary centered normals $(X,Z)-vq_0W$ supply the remaining
positive directions. Requiring separate positive pivots for the
relative normal and the scale difference would obscure this
completion.

This proves an exact leading-model identity on its stated bounded
subdomain. It does not discard the additional large source scales:
those are handled algebraically by the centered-packet reduction in
Section~\ref{sec:ui-full-parent-reduction}. The complete local lifted
Gram is proved in Section~\ref{sec:local-ui-certificate}; the outer
projective join remains a separate obligation.
